\section{The \Rome{} type system}
\label{sec:Rome-calculus}

This section declaratively defines the \Rome type system. 

\subsection{Syntax}
\begin{figure}  
\begin{gather*}
\begin{array}{r@{\hspace{5px}}l@{\qquad}r@{\hspace{5px}}l@{\qquad}r@{\hspace{5px}}l@{\qquad}r@{\hspace{5px}}l}
\text{Type variables} & \alpha \in \mathcal A & \text{Labels} & \ell \in \mathcal L
\end{array}
\\[5px]
\begin{doublesyntaxarray}
  \mcl{\text{Kinds}} & \kappa & ::= & \TypeK \mid \shade{\LabK \mid \RowK \kappa} \mid \kappa \to \kappa \\
  \mcl{\text{Predicates}} & \pi & ::= & \shade{\LeqP \rho \rho \mid \PlusP \rho \rho \rho} \\
  \text{Types} & \mcr{\Types \ni \phi, \tau, \rho, \xi} & ::= & \alpha \mid T \mid \shade{\pi \then \tau} \\
               &                              &     & \mid    & \forall \alpha\co\kappa. \tau \mid \lambda \alpha\co\kappa. \tau \mid \tau \, \tau \\
               &                              &     & \mid    & \shade{\RowIx i 0 m {\LabTy {\ell_i} {\tau_i}}\mid \ell} \\ 
               &                              &     & \mid    & \shade{{\Row {\LabTy {\xi} {\tau}}} \mid \Mapp{\phi}{\rho} \mid \rho \Compl \rho} \\
  \mcl{\text{Type constants}} & T & ::= & (\to) \mid  \mu \mid \shade{\KFam \Pi \kappa \mid \KFam \Sigma \kappa} \\\end{doublesyntaxarray}
\end{gather*}
\caption{Syntax}
\label{fig:syntax-types}
\end{figure}
The syntax of \Rome kinds, predicates, and types are given in \cref{fig:syntax-types}. We have highlighted the syntax
that extends system \Fome. 

% Two subtleties warrant explanation. 

% \begin{itemize}
% \item We distinguish row literals from \emph{row singletons} $\Row {\LabTy {\xi} {\tau}}$, which are of length one but permitted to have unreduced and/or variable label $\xi$. It is imperative that we distinguish row literals from row singletons, as permitting label variables in row literals permits an opportunity for row labels to overlap upon instantiation.
% \item We exclude the application of $\Pi$ and $\Sigma$ to any kind $\kappa$ that would yield a label-valued row, e.g. $\RowK \LabK$. This ensures that labels are canonically literal when in normal form. (Consider the nonsensical but otherwise permitted label $\Pi {\Row {\LabTy {\ell} {\ell}}}$, which would have kind $\LabK$.)
% \end{itemize}



\subsection{Kinding}

\begin{figure}
\fbox{$\KindJ \Delta \pi {\RowK \kappa}$}\, \fbox{$\KindJ \Delta \tau \kappa$}
{\small
\begin{gather*}
\ib{\irule[\ksrule{$\lesssim$}]
  {\KindJ \Delta {\rho_{i}} {\RowK \kappa}};
  {\KindJ \Delta {\LeqP {\rho_{1}} {\rho_{2}}} {\RowK \kappa}}}
\isp
\ib{\irule[\ksrule{$\cdot$}]
  {\KindJ \Delta {\rho_{i}} {\RowK \kappa}};
  {\KindJ \Delta {\PlusP {\rho_{1}} {\rho_{2}} {\rho_3}} {\RowK \kappa}}}
\\
\ib{\irule[\ksrule{lab}]
          {};
          {\KindJ \Delta \ell \LabK}}
\isp
\ib{\irule[\ksrule{$\then$}]
  {\KindJ \Delta \pi {\RowK \kappa}} 
  {\KindJ \Delta \tau \TypeK};
  {\KindJ \Delta {\pi \then \tau} {\TypeK}}}
\\
\ib{\irule[\ksrule{row}]
          {\KindJ \Delta {\ell_i} \LabK}
          {\KindJ \Delta {\tau_i} \kappa}
          {\forall i < j. \, \ell_i < \ell_j};
          {\KindJ \Delta {\RowIx i 0 n {\LabTy {\ell_i} {\tau_i}}} {\RowK \kappa}}}
\\
\rsp 
\ib{\irule[\krrule{rsing}]
          {\KindJ \Delta {\xi} \LabK}
          {\KindJ \Delta {\tau} \kappa};
          {\KindJ \Delta {\{\LabTy \xi \tau\}} {\RowK \kappa}}}          
\\
\ib{\irule[\krrule{map}]
          {\KindJ \Delta \phi {\kappa_1 \to \kappa_2}}
          {\KindJ \Delta \rho {\RowK {\kappa_1}}};
          {\KindJ \Delta {\Mapp \phi \rho } {\RowK{\kappa_2}}}}
\\
\ib{\irule[\krrule{compl}]
          {\KindJ \Delta {\rho_2} {\RowK\kappa}}
          {\KindJ \Delta {\rho_1} {\RowK\kappa}};
          {\KindJ \Delta {\rho_2 \Compl \rho_1} {\RowK\kappa}}}
\\[5pt]
\fbox{$T : \kappa$}
\\
\begin{aligned}
(\to) &: \TypeK \to \TypeK \to \TypeK \\
\mu &: (\TypeK \to \TypeK) \to \TypeK \\ 
\end{aligned}
\qquad\qquad
\begin{aligned}
\KFam \Pi \kappa &: \RowK \kappa \to \kappa \\
\KFam \Sigma \kappa &: \RowK \kappa \to \kappa 
\end{aligned}\\ 
\begin{aligned}
\text{where } & \kappa \notin \mathcal K = \{ \LabK, \RowK \kappa, \kappa' \to \kappa \mid  \kappa \in \mathcal K \} 
\end{aligned}
\end{gather*}}
\caption{Kinding}
\label{fig:kinding}
\end{figure}

\Cref{fig:kinding} describes the kind system of \Rome. RowLiteral (Type literals are) 


\subsection{Type Equivalence and Reduction}
\InlineOff{}
\begin{figure}
\begin{small}
\begin{gather*}
\fbox{$\TEqvJ \Delta \tau \tau \kappa$}
\\
{\footnotesize
\ib{\irule[\erule{$\beta$}]
          {\KindJ \Delta {(\lambda \alpha\co\kappa. \tau) \, \upsilon} {\kappa'}};
          {\TEqvJ \Delta {(\lambda \alpha\co\kappa. \tau)\,\upsilon} {\tau[\upsilon/\alpha]} {\kappa'}}}          
\isp 
\ib{\irule[\erule{$\eta$}]
  {\KindJ \Delta \phi {\kappa_1 \to \kappa_2}};
  {\TEqvJ \Delta {\phi} {\lambda \alpha. \phi\, \alpha} {\kappa_1 \to \kappa_2}}}}
\\
\ib{\irule[\errule{map}]
          {\KindJ \Delta \phi {\kappa_1 \to \kappa_2}}
          {\KindJ \Delta {\RowIx i 0 n {\LabTy {\ell_i} {\tau_i}}} {\RowK{\kappa_1}}};
          {\TEqvJ \Delta 
          {\Mapp {\phi} {\RowIx i 0 n {\LabTy {\ell_i} {\tau_i}}}} 
          {\RowIx i 0 n {\LabTy {\ell_i} {\phi\,\tau_i}}} 
          {\RowK{\kappa_2}}}}
\\
\ib{\irule[\errule{map$_\Labels$}]
          {\KindJ \Delta \phi {\kappa_1 \to \kappa_2}}
          {\KindJ \Delta {\Row {\LabTy {\xi} {\tau}}} {\RowK{\kappa_1}}};
          {\TEqvJ \Delta 
            {\Mapp {\phi} {\Row {\LabTy {\xi} {\tau}}}} 
            {\Row {\LabTy {\xi} {\phi\,\tau}}} 
            {\RowK{\kappa_2}}}}
\\
\ib{\irule[\errule{map$_\mathsf{id}$}]
          {\KindJ \Delta \rho {\RowK\kappa}};
          {\TEqvJ \Delta \rho {\Mapp {(\lambda \alpha. \alpha)} \rho} {\RowK\kappa}}}
\\
\begin{gathered}
\ib{\irule[\errule{map$_\circ$}]
          {\KindJ \Delta {\phi_1} {\kappa_2 \to \kappa_3}}
          {\KindJ \Delta {\phi_2} {\kappa_1 \to \kappa_2}}
          {\KindJ \Delta \rho {\RowK {\kappa_1}}};
          {\TEqvJ \Delta {\Mapp {\phi_1} {(\Mapp {\phi_2} \rho)}} {\Mapp {(\phi_1 \circ \phi_2)} \rho} {\kappa_3}}}
\\
\text{where $\phi_1 \circ \phi_2 = \lambda \alpha. \phi_1 \, (\phi_2 \, \alpha)$}
\end{gathered}
\\
\ib{\irule[\erule{$\setminus$}]
          {\KindJ \Delta {\RowIx i 0 n {\LabTy {\xi_i} {\tau_i}}} {\RowK \kappa}}
          {\KindJ \Delta {\RowIx j 0 m {\LabTy {\xi_j} {\tau_j}}} {\RowK \kappa}};
          % {\Delta \vdash \mathsf{subtract} \, \rho_1 \, \rho_2 \, \rho_3};
          {\TEqvJ \Delta {{\Row {\LabTy {\xi_i} {\tau_i}}} \Compl {\Row {\LabTy {\xi_j} {\tau_j}}}} {\Subtract \, {\Row {\LabTy {\xi_i} {\tau_i}}} \, {\Row{\LabTy {\xi_j} {\tau_j}}}} {\RowK \kappa}}}
\\
\ib{\irule[\errule{map$_{\setminus}$}]
          {\KindJ \Delta {\phi} {\kappa_1 \to \kappa_2}}
          {\KindJ \Delta {\rho_i} {\RowK {\kappa_1}}};
          {\TEqvJ \Delta {\Mapp {\phi} {(\rho_2 \setminus \rho_1)}} {\Mapp \phi {\rho_2} \setminus {\Mapp \phi {\rho_1}}} {\kappa_2} }}          
\\
\begin{gathered}
\ib{\irule[\errule{assoc$_\Xi$}]
          {\KindJ \Delta \rho {\RowK {\kappa_1 \to \kappa_2}}}
          {\KindJ \Delta \tau \kappa_1};
          {\TEqvJ \Delta {(\KFam \Xi {\kappa_1 \to \kappa_2} \rho) \, \tau} {\KFam \Xi {\kappa_2} (\rho \, \Flap \, \tau)} {\kappa_2}}!
          {\Xi \in \Set {\Pi, \Sigma}}}
\\
\text {where $\rho \, \Flap \, \tau = \Mapp {(\lambda f. f \, \tau)} \rho$}
\end{gathered}          
\\
\ib{\irule[\errule{$\Xi$}]
          {\KindJ \Delta \rho {\RowK {\RowK {\kappa}}}};
          {\TEqvJ \Delta {\KFam \Xi {\RowK \kappa} \, \rho} {\Mapp {{\KFam \Xi \kappa}} \rho} {\RowK {\kappa}}}!
          {\Xi \in \{ \Pi , \Sigma \}}}
\\[5pt]
\fbox{$\Subtract \, \rho \, \rho$}
\\
\begin{aligned}
  \Subtract \, \EmptyRow \, \rho &= \EmptyRow \\
  \Subtract \, \Row {\LabTy \ell \tau, \rho_2} \, \rho_1 &= 
    &\begin{cases}
      \Subtract \, \rho_2 \, \rho_1 &\text{if $\ell \in \rho_{1}$} \\
      \Row {\LabTy \ell \tau, \Subtract \, \rho_2 \, \rho_1} &\text{otherwise}
    \end{cases}
\end{aligned}
\end{gather*}
\end{small}
\caption{Type equivalence}
\label{fig:equivalence}
\end{figure}
\InlineOn{}

We define reduction on types $\tau \RedT \tau'$ by directing the type equivalence judgment $\TEqvJ \Delta \tau {\tau'} \kappa$ (\cref{fig:equivalence}) from left to right. The full equivalence relation is the congruence closure of \cref{fig:equivalence} plus rules for reflexivity, symmetry, and transitivity.

\subsection{Normal forms}
\label{sec:normal-forms}
\begin{figure}
\begin{gather*}
\begin{array}{r@{\hspace{7px}}l@{\qquad\qquad}r@{\hspace{7px}}l}
  \text{Type variables} & \alpha \in \mathcal A &
  \text{Labels} & \ell \in \mathcal L
\end{array} \\
\begin{doublesyntaxarray}
  \mcl{\text{Ground Kinds}}  & \gamma   & ::= & \TypeK \mid \LabK \\
  \mcl{\text{Kinds}}         & \kappa    & ::= & \shade{\gamma} \mid \kappa \to \kappa \mid  \RowK \kappa \\
  \mcl{\text{Predicates}} & {\Normal \pi} & ::= & \LeqP {\Normal \tau} {\Normal \tau} \\
  \mcl{\text{Row Literals}}   & \NormalRows \ni \Normal \rho    & ::= & \shade{\RowIx i 0 m {\LabTy {\ell_i} {\Normal {\tau_i}}}} \\
  \mcl{\text{Neutral Types}} & n    & ::= & \alpha \mid n \, {\Normal \tau}  \\
  \mcl{\text{Normal Types}}  & \NormalTypes \ni \Normal \tau, \Normal \phi & ::= & n \mid \shade{\Mapp {\hat{\phi}} {n} \mid \Normal{\rho}} \mid \Normal{\pi} \then \Normal{\tau} \\
                             &       &     & \mid &  \forall \alpha\co\kappa. \Normal{\tau} \mid \lambda \alpha\co\kappa. \Normal{\tau}  \\ 
                             &       &     & \mid & \Normal \tau \to \Normal \tau \mid \mu\, \Normal \tau \\
                             &       &     & \mid & \shade{\{\LabTy n {\Normal \tau}\}} \mid \ell \mid \shade{{\Normal \tau} \Compl {\Normal \tau}}  \\
                             &       &     & \mid & \shade{{\KFam \Pi \TypeK} {\Normal \tau} \mid {\KFam \Sigma \TypeK} {\Normal \tau}} 
\end{doublesyntaxarray}
\end{gather*}
\begin{small}
\begin{gather*}
\fbox{$\KindJNF \Delta {\Normal \tau} {\kappa}$} \, \fbox{$\KindJNE \Delta {n} {\kappa}$} \\
\ib{
  \irule[\kruleNF{ne}]
    {\KindJNE \Delta n \shade{\gamma}};
    {\KindJNF \Delta n \shade{\gamma}}}
\rsp
\ib{
  \irule[\kruleNF{$\triangleright$}]
    {\shade{\KindJNE \Delta {n} {\LabK}}}{\KindJNF \Delta {\Normal \tau} \kappa};
    {\KindJNF \Delta {\LabTy n {\Normal \tau}}{\RowK \kappa}}}
\\
\ib{
  \irule[\kruleNF{$\Compl$}]
    {\KindJNF \Delta {\Normal {\tau_i}} {\RowK \kappa}}
    {\shade{\Normal{\tau_1} \notin \NormalRows \, \text{or}\, \Normal{\tau_2} \notin \NormalRows}};
    {\KindJNF \Delta {\Normal{\tau_2} \Compl \Normal{\tau_1}} {\RowK \kappa}}}
\end{gather*}
\end{small}
\caption{Normal type forms}
\label{fig:type-normalization}
\end{figure}

The syntax and kinding of normal and neutral types is given in \cref{fig:type-normalization}. % We have highlighted some subtle aspects of the syntax and kinding for discussion:

% \begin{itemize}
% \item we distinguish ground kinds $\gamma$ from functional and row kinds, as neutral types may only be promoted to normal type at ground kind (rule \kruleNF{ne});
% \item neutral types at row kind are expanded to an inert map by the identity function (e.g., to $\Mapp {(\lambda x. x)} {n}$);
% \item the syntax $\{\LabTy n \Normal \tau\}$ separates singleton rows (with neutral label $n$) from row literals $\Normal \rho$ (with literal labels $\ell_i$); rule \kruleNF{$\triangleright$} ensures that $n$ is a well-kinded neutral label;
% \item we restrict the formation of $\Pi$ and $\Sigma$ types to just kind $\TypeK$, as this enforces arrow-kinded types are canonically $\lambda$-bound; and
% \item  the complement of two row literals should compute to another row literal; we thus require in the kinding of normal row complements \kruleNF{$\Compl$} that one (or both) rows are not literal.
% \end{itemize}
